\documentclass[tikz,border=10pt]{standalone}
\usepackage[no-math]{fontspec}
\usepackage[brazil]{babel}
\usepackage{amsmath,amssymb}
\IfFileExists{C:/Windows/Fonts/Lato_Regular.ttf}{%
  \setmainfont{Lato}[
    Path = C:/Windows/Fonts/,
    UprightFont = Lato_Regular.ttf,
    BoldFont = Lato_Bold.ttf,
    ItalicFont = Lato_Italic.ttf,
    BoldItalicFont = Lato_BoldItalic.ttf
  ]
  \setsansfont{Lato}[
    Path = C:/Windows/Fonts/,
    UprightFont = Lato_Regular.ttf,
    BoldFont = Lato_Bold.ttf,
    ItalicFont = Lato_Italic.ttf,
    BoldItalicFont = Lato_BoldItalic.ttf
  ]
}{%
  \RequirePackage[default]{lato}
}
\renewcommand{\familydefault}{\sfdefault}

\usepackage{xcolor}
\usepackage{tikz}
\usepackage{pgfplots}
\pgfplotsset{compat=1.18}
\usetikzlibrary{arrows.meta,positioning,shapes.geometric,fit,backgrounds,calc,decorations.pathreplacing,patterns}

% Paleta institucional UEFS / Bioengenharia de Solos
\definecolor{bsNavy}{HTML}{1F4268}
\definecolor{bsNavyDark}{HTML}{152E4A}
\definecolor{bsGold}{HTML}{B38A3D}
\definecolor{bsGoldDark}{HTML}{8C6927}
\definecolor{bsTeal}{HTML}{0E7C7B}
\definecolor{bsForest}{HTML}{2D6A4F}
\definecolor{bsAlert}{HTML}{9E2A2B}
\definecolor{bsEarth}{HTML}{8C5A3C}
\definecolor{bsWater}{HTML}{2E86AB}
\definecolor{bsInk}{HTML}{26313D}
\definecolor{bsMuted}{HTML}{5B6573}

\definecolor{bgNavy}{HTML}{EDF2F7}
\definecolor{bgGold}{HTML}{F8F3E9}
\definecolor{bgTeal}{HTML}{EBF6F5}
\definecolor{bgForest}{HTML}{EEF5F0}
\definecolor{bgAlert}{HTML}{FDF2F2}

\tikzset{
  >=Stealth,
  bsBox/.style={
    draw=#1, fill=white, rounded corners=5pt, line width=1.15pt,
    inner sep=7pt, align=center, font=\footnotesize\sffamily, text=bsInk
  },
  bsTint/.style={
    draw=#1, fill=#1!8!white, rounded corners=5pt, line width=1.15pt,
    inner sep=7pt, align=center, font=\footnotesize\sffamily, text=bsInk
  },
  bsBadge/.style={
    fill=#1, text=white, rounded corners=9pt, inner xsep=8pt, inner ysep=3pt,
    font=\scriptsize\bfseries\sffamily
  },
  bsTitle/.style={font=\normalsize\bfseries\sffamily, text=bsNavy, align=center},
  bsSub/.style={font=\footnotesize\sffamily, text=bsMuted, align=center},
  bsArrow/.style={->, >=Stealth, line width=1.35pt},
  bsDArrow/.style={<->, >=Stealth, line width=1.35pt}
}


\begin{document}
\begin{tikzpicture}[font=\sffamily]

% ============================================================
% TÍTULO GERAL
% ============================================================
\node[bsTitle, anchor=west] at (0, 8.35)
  {Hidráulica de Encostas e Dimensionamento do Espaçamento entre Cordões};
\node[bsSub, anchor=west] at (0, 7.90)
  {Segmentação do comprimento de rampa, terraçamento progressivo e critérios de projeto};

% ============================================================
% PAINEL ESQUERDO: ESQUEMA HIDRÁULICO NA ENCOSTA
% ============================================================
\begin{scope}[shift={(0,0)}]
  \draw[bsWater!60, rounded corners=6pt, line width=1.1pt, fill=white]
    (0,-0.55) rectangle (10.10, 7.45);
  \node[bsBadge=bsWater, rounded corners=5pt, anchor=north west] at (0.30, 7.20) {HIDRÁULICA DE ENCOSTA};

  % Terreno original e patamares de sedimentação
  \fill[bsEarth!16]
    (0.45, 0.70) -- (0.45, 5.55) -- (8.85, 1.75) -- (8.85, 0.70) -- cycle;

  % Cunhas de terraço progressivo (sedimentos depositados a montante de cada cordão)
  \fill[bsGold!40] (1.15, 5.23) -- (3.35, 4.24) -- (3.35, 4.78) -- cycle;
  \fill[bsGold!40] (3.45, 4.19) -- (6.15, 2.97) -- (6.15, 3.55) -- cycle;

  % Perfil original da encosta
  \draw[bsEarth!75!black, dashed, line width=1.0pt] (0.45, 5.55) -- (8.85, 1.75);
  % Superfície escalonada (terraço progressivo)
  \draw[bsGoldDark, line width=1.3pt] (0.45, 5.55) -- (1.15, 5.23) -- (3.35, 4.78);
  \draw[bsGoldDark, line width=1.3pt] (3.45, 4.19) -- (6.15, 3.55);
  \draw[bsEarth, line width=1.2pt] (6.25, 2.92) -- (8.85, 1.75);

  % Cordão 1 (superior)
  \foreach \dx/\h in {-0.12/0.72, -0.04/0.86, 0.04/0.92, 0.12/0.82, 0.20/0.68}{
    \draw[bsForest, line width=1.4pt, line cap=round] (3.40+\dx, 4.22) -- (3.40+\dx*1.4, 4.22+\h);
  }
  \draw[bsEarth!80!black, line width=0.8pt] (3.40, 4.20) -- (3.25, 3.52);
  \draw[bsEarth!80!black, line width=0.8pt] (3.46, 4.20) -- (3.58, 3.48);

  % Cordão 2 (inferior)
  \foreach \dx/\h in {-0.12/0.72, -0.04/0.86, 0.04/0.92, 0.12/0.82, 0.20/0.68}{
    \draw[bsForest, line width=1.4pt, line cap=round] (6.20+\dx, 2.95) -- (6.20+\dx*1.4, 2.95+\h);
  }
  \draw[bsEarth!80!black, line width=0.8pt] (6.20, 2.93) -- (6.05, 2.25);
  \draw[bsEarth!80!black, line width=0.8pt] (6.26, 2.93) -- (6.38, 2.22);

  % Setas de escoamento superficial desacelerado e drenagem limpa (bem separadas dos rótulos)
  \draw[bsArrow, bsWater, line width=1.6pt] (0.75, 5.95) -- (2.95, 5.15);
  \node[font=\tiny\bfseries\sffamily, text=bsWater, anchor=south west] at (0.70, 6.02)
    {Desaceleração do escoamento};

  \draw[bsArrow, bsWater!75!black, line width=1.15pt, dashed] (3.85, 4.52) -- (5.75, 3.92);
  \draw[bsArrow, bsWater!75!black, line width=1.15pt, dashed] (6.70, 3.22) -- (8.65, 2.34);
  \node[font=\tiny\bfseries\sffamily, text=bsWater!85!black, align=center, anchor=south] at (7.85, 2.95)
    {Drenagem limpa\\a jusante};

  % Rótulos das cunhas e cordões sem colisões
  \node[font=\tiny\bfseries\sffamily, text=bsGoldDark, align=center, anchor=north east] (lTer) at (2.85, 4.05)
    {Terraço\\Progressivo};
  \draw[bsArrow, bsGoldDark, line width=0.75pt] (lTer.north) -- (2.55, 4.78);

  \node[font=\tiny\bfseries\sffamily, text=bsGoldDark, align=center, anchor=north] (lCun) at (4.90, 3.25)
    {Cunha de sedimentos};
  \draw[bsArrow, bsGoldDark, line width=0.75pt] (lCun.north) -- (5.05, 3.68);

  \node[font=\tiny\bfseries\sffamily, text=bsForest, align=left, anchor=west] at (3.75, 5.18)
    {Cordão Vegetativo};

  % Cotas EV (Espaçamento Vertical) e EH (Espaçamento Horizontal)
  \draw[bsNavy, dotted, line width=0.8pt] (3.40, 4.22) -- (6.50, 4.22);
  \draw[bsNavy, dotted, line width=0.8pt] (6.20, 2.95) -- (6.50, 2.95);
  \draw[bsDArrow, bsNavy, line width=1.05pt] (6.45, 2.95) -- (6.45, 4.22)
    node[midway, right=2pt, font=\scriptsize\bfseries\sffamily, text=bsNavy] {$EV$};

  \draw[bsNavy!55, dotted, line width=0.8pt] (3.40, 4.22) -- (3.40, 1.15);
  \draw[bsNavy!55, dotted, line width=0.8pt] (6.20, 2.95) -- (6.20, 1.15);
  \draw[bsDArrow, bsNavy, line width=1.05pt] (3.40, 1.15) -- (6.20, 1.15)
    node[midway, below=3pt, font=\scriptsize\bfseries\sffamily, text=bsNavy] {$EH$ (Espaçamento Horizontal)};

  % Direção de implantação de baixo para cima (mais à direita, livre do rótulo de drenagem)
  \draw[bsArrow, bsForest, line width=1.35pt] (9.35, 1.85) -- (9.35, 5.45)
    node[midway, sloped, above=3pt, font=\tiny\bfseries\sffamily, text=bsForest] {Implantação: baixo $\rightarrow$ cima};
\end{scope}

% ============================================================
% PAINEL DIREITO: REGRAS DE DIMENSIONAMENTO (bsNavy)
% ============================================================
\begin{scope}[shift={(10.50,0)}]
  \draw[bsNavy, rounded corners=6pt, line width=1.2pt, fill=bgNavy]
    (0,-0.55) rectangle (5.65, 7.45);
  \node[bsBadge=bsNavy, rounded corners=5pt, anchor=north west] at (0.30, 7.20) {CRITÉRIOS DE PROJETO};

  \node[bsBox=bsNavy, anchor=north west, text width=4.75cm, align=left] (r1) at (0.30, 6.58)
    {\textbf{\color{bsNavy}1. Variáveis Condicionantes}\\
     $\bullet$ Declividade do terreno ($S\%$)\\
     $\bullet$ Erodibilidade do solo ($K$)\\
     $\bullet$ Erosividade da chuva ($R$)};

  \node[bsBox=bsWater, anchor=north west, text width=4.75cm, align=left] (r2) at ($(r1.south west)+(0,-0.45)$)
    {\textbf{\color{bsWater}2. Cálculo do Espaçamento}\\
     Determinação do Espaçamento Vertical ($EV$) e Horizontal ($EH = EV / \tan\theta$), limitando a energia do escoamento entre faixas.};

  \node[bsBox=bsForest, anchor=north west, text width=4.75cm, align=left] (r3) at ($(r2.south west)+(0,-0.45)$)
    {\textbf{\color{bsForest}3. Ordem de Execução}\\
     Implantação \textbf{de baixo para cima} rigorosamente em \textbf{curvas de nível}, garantindo ancoragem antes do aporte a montante.};

  \draw[bsArrow, bsNavy] (r1.south) -- (r2.north);
  \draw[bsArrow, bsNavy] (r2.south) -- (r3.north);
\end{scope}

\end{tikzpicture}
\end{document}
